\newpage
\section*{РЕЗЮМЕ СТАРТАП-ПРОЕКТА}
\addcontentsline{toc}{section}{РЕЗЮМЕ СТАРТАП-ПРОЕКТА}


Название: <<Бизнес-проект <<Приложение для составления тренировочного плана и отслеживания прогресса.>> Клиентская часть>>. Проект представляет собой программно-информационную систему, предназначенную для организации самостоятельных занятий физической культурой: составления индивидуального плана нагрузки, учёта выполненных упражнений, контроля питания и обмена результатами между пользователями. Система обеспечивает регистрацию и вход в личный кабинет, ведение профиля с учётом целей и антропометрических данных, формирование и хранение тренировочных планов, фиксацию подходов и анализ прогресса, ведение дневника питания, работу социальной ленты, а также синхронизацию данных между браузером пользователя и сервером с защитой персональной информации.

Ключевые возможности системы включают:
\begin{enumerate}
	\item Составление и редактирование индивидуальных планов тренировок с учётом целей, уровня подготовки и предпочтений пользователя.
	
	\item Ведение журнала выполненных упражнений с фиксацией повторений, рабочего веса и нагрузки по группам мышц, а также наглядное отображение динамики прогресса и сводной статистики.
	
	\item Учёт питания и пользовательских настроек (цели по калориям и макронутриентам, собственные продукты) в едином веб-интерфейсе.
	
	\item Социальную ленту: публикации о тренировках и питании, комментарии, лайки и реакции для поддержания мотивации и обмена опытом.
	
	\item Интуитивно понятный и адаптивный графический интерфейс веб-клиента, поддерживающий работу на различных устройствах без установки отдельного мобильного приложения.
\end{enumerate}

Созданная система оптимизирует процесс самостоятельных тренировок, повышая наглядность контроля нагрузки и прогресса и позволяя пользователю объединить планирование занятий, учёт питания и социальное взаимодействие без переключения между несколькими узкоспециализированными сервисами.

\emph{Актуальность} проекта обусловлена ростом интереса к здоровому образу жизни и цифровым инструментам самоконтроля. Большинство популярных решений закрывают лишь отдельные задачи -- учёт силовых тренировок, дневник питания или социальную мотивацию -- и часто требуют подписки либо привязаны к конкретной платформе. Пользователю, занимающемуся без постоянного тренера, необходим доступный веб-сервис, который помогает организовать занятия, видеть результат и сохранять данные при работе с разных устройств, что снижает затраты времени на подбор разрозненных приложений.

\emph{Целью} разработки является создание программно-информационной системы <<Fitness Life>>, способной обеспечить:
\begin{itemize}
	\item персональное планирование тренировок и учёт выполненных занятий;
	
	\item отслеживание прогресса и наглядную аналитику физической нагрузки;
	
	\item ведение дневника питания в составе единого приложения;
	
	\item социальное взаимодействие пользователей для поддержания мотивации;
	
	\item надёжное хранение данных и доступ к личному кабинету с различных устройств через браузер.
\end{itemize}

\emph{Отличительные особенности} разработанной системы по сравнению с существующими решениями:
\begin{enumerate}
	\item Современный и интуитивно понятный интерфейс, адаптированный под сценарии планирования тренировок, учёта питания и работы с лентой.
	
	\item Объединение тренировочного плана, журнала занятий, учёта питания и социальной ленты в одном web-приложении.
	
	\item Гибкая настройка профиля и целей пользователя, облачная синхронизация данных между сеансами работы и устройствами.
	
	\item Возможность масштабирования и дальнейшей интеграции с мобильными платформами, носимыми устройствами и внешними фитнес-сервисами.
\end{enumerate}

В системе реализованы задачи по отображению экранов регистрации и входа, ведению профиля, формированию и хранению тренировочных планов на клиенте, фиксации выполненных упражнений и визуализации статистики прогресса, учёту питания, работе социальной ленты и синхронизации данных браузера с REST API серверной части, а также обеспечению модульной структуры клиентского приложения с использованием современных веб-технологий (HTML, CSS, JavaScript, Chart.js, localStorage).

Идея создания проекта возникла в 2025 году на фоне возросшего спроса на доступные цифровые инструменты для самостоятельных тренировок с облачным хранением прогресса, а разработка была начата в начале 2026 года. Завершение проекта планируется к середине 2026 года.

\emph{Бизнес-цели:}
\begin{enumerate}
	\item Расширение функциональных возможностей за счёт интеграции с трекерами активности, каталогами упражнений и внешними сервисами здоровья.
	
	\item Масштабирование системы для роста пользовательской базы с учётом различных целей занятий (силовые тренировки, похудение, поддержание формы).
	
	\item Добавление дополнительных функций: напоминания о тренировках, визуализация серий занятий, расширенные шаблоны программ и настройка интерфейса для повышения вовлечённости.
	
	\item Повышение регулярности занятий и удержание аудитории за счёт наглядного прогресса, социальной поддержки и модели freemium (базовый функционал бесплатно, расширенная аналитика -- по подписке).
\end{enumerate}
